\par\Needspace{6\baselineskip}\medskip
\begin{center}\hypertarget{native-235-article170}{}\pdfbookmark[2]{170.}{article-170}170.\end{center}
In hypothesi prima quidem temporibus non correctis $t,t',t'',t'''$ uti oportebit, quum distantias a terra computare nondum liceat: simulac vero valores approximati quantitatum $x',x''$ innotuerunt, illas distantias quoque proxime determinare poterimus. Attamen quum formulae pro $\rho$ et $\rho'''$ hic paullo complicatiores evadant, computum correctionis temporum eousque differre conveniet, ubi distantiarum valores satis praecisi evaserunt, ne calculo repetito opus \SourceLeaf{236}sit. Quamobrem e re erit, hanc operationem iis valoribus quantitatum $x',x''$ superstruere, ad quas hypothesis penultima produxit, ita ut ultima demum hypothesis a valoribus correctis temporum atque quantitatum $P',P'',Q',Q''$ proficiscatur. Ecce formulas, ad hunc finem in usum vocandas:
\[
\tag*{VII.}
\rho'=x'-R'\cos\delta'
\]
\[
\tag*{VIII.}
\rho''=x''-R''\cos\delta''
\]
\[
\tag*{IX.}
\begin{aligned}
\rho\cos\beta={}&-R\cos B\cos(a-l)\\
&+\frac{1+P'}{P'\left(1+\dfrac{Q'}{r'^3}\right)}\bigl(\rho'\cos\beta'\cos(a'-a)+R'\cos B'\cos(l'-a)\bigr)\\
&-\frac{1}{P'}\bigl(\rho''\cos\beta''\cos(a''-a)+R''\cos B''\cos(l''-a)\bigr)
\end{aligned}
\]
\[
\tag*{X.}
\begin{aligned}
\rho\sin\beta={}&-R\sin B\\
&+\frac{1+P'}{P'\left(1+\dfrac{Q'}{r'^3}\right)}\bigl(\rho'\sin\beta'+R'\sin B'\bigr)\\
&-\frac{1}{P'}\bigl(\rho''\sin\beta''+R''\sin B''\bigr)
\end{aligned}
\]
\[
\tag*{XI.}
\begin{aligned}
\rho'''\cos\beta'''={}&-R'''\cos B'''\cos(a'''-l''')\\
&+\frac{1+P''}{P''\left(1+\dfrac{Q''}{r''^3}\right)}\bigl(\rho''\cos\beta''\cos(a'''-a'')\\
&\qquad{}+R''\cos B''\cos(a'''-l'')\bigr)\\
&-\frac{1}{P''}\bigl(\rho'\cos\beta'\cos(a'''-a')+R'\cos B'\cos(a'''-l')\bigr)
\end{aligned}
\]
\[
\tag*{XII.}
\begin{aligned}
\rho'''\sin\beta'''={}&-R'''\sin B'''\\
&+\frac{1+P''}{P''\left(1+\dfrac{Q''}{r''^3}\right)}\bigl(\rho''\sin\beta''+R''\sin B''\bigr)\\
&-\frac{1}{P''}\bigl(\rho'\sin\beta'+R'\sin B'\bigr).
\end{aligned}
\]
Formulae IX—XII nullo negotio ex aequationibus 1, 2, 3 art. 112 derivantur, si modo characteres illic adhibiti in eos quibus hic utimur rite convertuntur. Manifesto formulae multo simpliciores evadunt, si $B,B',B''$ evanescunt. E combinatione formularum IX et X non modo $\rho$ sed etiam $\beta$, et perinde ex XI et XII praeter $\rho'''$ etiam $\beta'''$ demanat: valores harum latitudinum cum observatis (calculum non ingredientibus), siquidem datae sunt, comparati ostendent, quonam praecisionis gradu latitudines extremae per elementa sex reliquis datis adaptata repraesentari possint.
